\section{Gluing the vector-field jets}\label{jets:section}

We use the curve and sections of Proposition~\ref{prop:curve-data},
with the normalization supplied by Proposition~\ref{prop:residue-data}.
Thus $L^4\simeq K_C$, $\deg L=10$, and
$r=(r_1,r_2)=\xi(u_1,u_2)$ is a basis of $H^0(C,L)$ with simple common
zero divisor $Z=q_1+\cdots+q_6$.
For rows $X,Y$ put $X\wedge Y=X_1Y_2-X_2Y_1$.
In the compatible coordinates at $q_j$, the preceding proposition gives
\begin{equation}\label{jets:residue-input}
 \begin{gathered}
 m_j=\frac{c_j}{z},\qquad R_j=m_jr,\qquad p=P-m_jr',\qquad
 p(0)\wedge R_j(0)=1,\\
 \sum_{j=1}^6\bigl(W_j-1\bigr)=-9\deg L,
 \qquad W_j=R_j(0)\wedge R_j'(0).
 \end{gathered}
\end{equation}
Here $P$ is a meromorphic row of sections of $L^2$, holomorphic off $Z$,
and $p$ and $R_j$ are holomorphic on the disk in question.
The projective directions of $r'(q_j)$ are not all equal.
Primes always denote differentiation in that disk's base coordinate.
Our task is to choose finitely many neighborhood transitions and compatible
vector-field coefficients whose determinant has the required leading term.
We first solve the exceptional low-order obstruction, then extend the fields
and their determinant together. Proposition~\ref{prop:jets} gives the finite
jet needed for holomorphic lifting.

\subsection{Coordinates, transitions, and obstruction spaces}

On the total space of $L^{-1}$ choose a base coordinate $z$ and a
fiber-linear equation $t$ of the zero section such that $t^4$ corresponds
to $dz$ under $L^4\simeq K_C$.  Two such choices satisfy
\[
 \widetilde z=\psi(z),\qquad \widetilde t=g(z)t,
 \qquad g^4=\psi'.
\]
Put $E_t=t\partial_t$.  The meromorphic two-form
\begin{equation}\label{jets:volume}
 \sigma=t^{-4}\frac{dt}{t}\wedge dz=t^{-5}dt\wedge dz
\end{equation}
is invariant under these split changes of coordinates.
We seek rows of coefficients $A_k,B_k$ defining
\begin{equation}\label{jets:fields}
 \begin{gathered}
 v_i=t\bigl(A_i(t,z)E_t+B_i(t,z)\partial_z\bigr),\\
 A=\sum_{k\geq0}A_kt^k,\quad B=\sum_{k\geq0}B_kt^k,
 \qquad A_0=r,\quad B_0=0.
 \end{gathered}
\end{equation}
The subscripts $k$ label orders; subscripts $i,j\in\{1,2\}$ label
entries of a row.  In particular,
\begin{equation}\label{jets:determinant}
 \sigma(v_1,v_2)=t^{-2}(A\wedge B),\qquad
 v_1\wedge v_2=t^2(A\wedge B)E_t\wedge\partial_z.
\end{equation}

For a term $t^{k+1}(A_kE_t+B_k\partial_z)$, a split coordinate change
gives, with $a=(\log g)'$,
\begin{equation}\label{jets:split-coefficients}
 \widetilde A_k=g^{-(k+1)}(A_k+aB_k),\qquad
 \widetilde B_k=g^{-(k+1)}\psi' B_k=g^{3-k}B_k.
\end{equation}
Consequently the horizontal coefficient has type $L^{k-3}$, and the
difference of two vertical lifts of a fixed horizontal coefficient has
type $L^{k+1}$.  More precisely, the sheaf of homogeneous terms of this
order is an extension
\begin{equation}\label{jets:homogeneous-extension}
 0\longrightarrow L^{k+1}\longrightarrow\mathcal E_k
 \longrightarrow L^{k-3}\longrightarrow0.
\end{equation}
This also specifies how rows written in local coordinates are interpreted
on a chart without a global coordinate or frame.

Take the cover consisting of $U_0=C\setminus Z$ and six mutually
disjoint disks $U_j$ about the $q_j$. The positive-degree divisor $Z$
has a very ample multiple; the corresponding projective embedding
realizes $U_0$ as a closed algebraic subvariety of the affine complement
of a hyperplane. Its analytification is therefore Stein. The disks and
annular intersections are Stein as well.  Thus horizontal sections admit holomorphic vertical lifts
on each chart, by \eqref{jets:homogeneous-extension} and Cartan's
Theorem~B~\cite[\S4, Theorem~B]{Cartan1952}.  Cousin problems on this cover compute the indicated first
cohomology groups.  Equivalently, one may shrink to a cover with annular
intersections.  For an overlap differential we use the contour residue
$\frac1{2\pi i}\int\omega$ around the missing point, oriented positively
in the disk coordinate.  The Serre-duality obstruction to a splitting
is the sum of these residues after multiplication by a dual section.
This formulation allows overlap functions holomorphic on annuli; they
need not be meromorphic at the omitted centers.

We modify the disk-to-outside identifications by flows of the vector
fields
\begin{equation}\label{jets:generator}
 D=t^k(bE_t+h\partial_z),\qquad k>0,
\end{equation}
where $b,h$ are holomorphic on the overlap in the disk coordinate.
Use the time-one flow of $-D$, so that its action on vector fields is
$\exp(\operatorname{ad}D)$.  Successive modifications are applied to
rows already expressed on the outside; all formulas below are relative
to the original split comparison.  Every generator vanishes on the
zero section and changes its equation only by $t\mapsto t+O(t^2)$,
so the zero section and its conormal bundle are preserved.

\begin{lemma}\label{jets:bracket}
For $D$ as in \eqref{jets:generator} and
$V=t^{\ell+1}(AE_t+B\partial_z)$ one has
\begin{align}
 [D,V]=t^{k+\ell+1}\bigl(&((\ell+1-k)bA+hA'-Bb')E_t\notag\\
 &+((\ell+1)bB+hB'-khA-Bh')\partial_z\bigr).
 \label{jets:general-bracket}
\end{align}
Moreover,
\begin{equation}\label{jets:divergence}
 \mathcal L_D\sigma=t^k\bigl((k-4)b+h'\bigr)\sigma.
\end{equation}
In particular, for $k\ne4$ the choice $b=h'/(4-k)$ preserves $\sigma$
exactly.  The first effect of an order-$k$ modification on
\eqref{jets:fields} is
\begin{equation}\label{jets:leading-action}
 \Delta A_k=(1-k)br+hr',\qquad \Delta B_k=-khr.
\end{equation}
\end{lemma}

\begin{proof}
Use $E_t(t^a)=at^a$ and $[E_t,\partial_z]=0$ in the two components of
the bracket.  For the divergence, the coefficient of $dt\wedge dz$ in
$\sigma$ is $t^{-5}$, so
\[
 \operatorname{div}_{\sigma}D
 =t^5\left(\partial_t(t^{k-4}b)
              +\partial_z(t^{k-5}h)\right)
 =t^k((k-4)b+h').
\]
Setting $\ell=0$, $A=r$, $B=0$ gives the last formula.
\end{proof}

All calculations can initially be made with finite Taylor series.
They nevertheless prescribe actual transition germs: on a compact
subannulus every $D$ in \eqref{jets:generator} is holomorphic near its
zero section and vanishes there, so its time-one flow and inverse exist
after shrinking in the transverse direction.  A finite composition has
the computed Taylor series.  There are no intersections between
different disk charts, hence no additional triple-overlap equations.
The passage from these germs to a Hausdorff surface neighborhood will
be made in the next section.

\subsection{The first two orders}

At order one use $h=m=m_j$, $b=m'/3$ on the $j$th overlap, and choose
\begin{equation}\label{jets:first-order}
\begin{array}{c|cc}
 &A_1&B_1\\ \hline
 U_j&p&mr\\
 U_0&P&0.
\end{array}
\end{equation}
By \eqref{jets:leading-action}, the first transition adds
$(mr',-mr)$ to the inside row, giving exactly $(P,0)$.
The degree-one coefficient of $A\wedge B$ is zero on both charts.

At order two write $U=B_2$.  To give the determinant its required
leading coefficient, solve on each chart
\begin{equation}\label{jets:second-horizontal}
 r\wedge U=1-A_1\wedge B_1.
\end{equation}
On a disk write $r=z\rho$, with $\rho$ nowhere zero as a row after
shrinking.  By \eqref{jets:residue-input}, the right side of
\eqref{jets:second-horizontal} vanishes at $z=0$, so division by $z$
reduces it to a surjective wedge map with $\rho$.
On $U_0$ the row $r$ has no common zero, and the corresponding
surjection of coherent sheaves is surjective on sections by
Steinness.  Thus the required holomorphic horizontal rows exist.
After the first transition, their horizontal difference on an overlap
wedges to zero with $r$: the determinant target is compatible because
that transition preserves $\sigma$.
The difference is therefore a unique holomorphic multiple of $r$.
An order-two transition with $b=h'/2$ removes it, since its horizontal
effect is $-2hr$.  Keep the chosen outside row $U$ fixed throughout this
matching.

It remains to match the vertical lifts $A_2$.  The obstruction belongs
to $H^1(C,L^3)^2$.  Its four Serre-duality tests use $r_1,r_2$, a basis
of $H^0(C,K_C\otimes L^{-3})=H^0(C,L)$, against the two entries of
the vertical difference.  The following calculation evaluates all four.

\begin{lemma}\label{jets:second-obstruction}
The order-two vertical obstruction vanishes.
\end{lemma}

\begin{proof}
Write $D^{(2)}=A_2$; this notation denotes the vertical row, not a
transition generator.  For $i,j\in\{1,2\}$ set
\begin{equation}\label{jets:psi}
 \Psi_{ij}=\left(r_iD_j^{(2)}+\frac14r_i'U_j
                          +\frac34r_j'U_i\right)dz.
\end{equation}
Once the horizontal rows have matched, the difference of these
expressions for the transported inside lift and the outside lift is
exactly $r_i$ times the $j$th vertical difference, multiplied by $dz$.
The derivative terms therefore change neither the obstruction nor its
four tests.

First compute the total for the outside lift.  Put
$\epsilon_{12}=1$, $\epsilon_{21}=-1$, and
$\epsilon_{11}=\epsilon_{22}=0$.
Under a split change with $a=(\log g)'$, formula
\eqref{jets:split-coefficients} gives
\[
 \widetilde r_i=g^{-1}r_i,\qquad
 \partial_{\widetilde z}\widetilde r_i=g^{-5}(r_i'-ar_i),\qquad
 \widetilde U_i=gU_i,\qquad
 \widetilde D_i^{(2)}=g^{-3}(D_i^{(2)}+aU_i).
\]
Substituting in \eqref{jets:psi}, and expressing both differentials in
the old coordinate, yields
\begin{equation}\label{jets:psi-transformation}
 \widetilde\Psi_{ij}-\Psi_{ij}
 =\frac34(r_iU_j-r_jU_i)d\log g
 =\frac34\epsilon_{ij}\,d\log g
 \quad\hbox{on }U_0.
\end{equation}
The last equality uses $r\wedge U=1$ on the outside.
Use the section $r_1=\xi u_1$ of $L$.  Its coefficient transforms by
$\widetilde r_1=g^{-1}r_1$, so
$\Psi_{ij}+\frac34\epsilon_{ij}d\log r_1$ is an intrinsic differential
on $U_0$.  There it has precisely the four possible simple poles at
the reduced zero fiber of $u_1$, each of residue
$\frac34\epsilon_{ij}$.  The residue theorem on the complement of
small disks therefore gives total contour residue $-3\epsilon_{ij}$
around $Z$ for this intrinsic differential.  At each of the six
points of $Z$, the section $r_1$ has a simple zero.  Subtracting these
six logarithmic residues yields
\begin{equation}\label{jets:outside-residue}
 \sum_{q\in Z}\Res_q\Psi_{ij}^{\mathrm{out}}
 =-\frac34\epsilon_{ij}\deg L.
\end{equation}
Only contour integrals at $Z$ were used, so this argument does not
assume meromorphic extension of the outside lifts.

Next compute the transported inside total.  Before transport,
$\Psi_{ij}$ is holomorphic on the disk.  The order-two transition
changes its rows by
\[
 \delta D^{(2)}=hr'-\frac12h'r,\qquad \delta U=-2hr.
\]
In \eqref{jets:psi} this is the exact differential
\begin{equation}\label{jets:psi-exact}
 \delta\Psi_{ij}=-\frac12d(hr_ir_j),
\end{equation}
so it has zero contour residue.
It suffices to calculate the effect of the order-one transition,
including its second exponential term.

The bracket of $t(m'E_t/3+m\partial_z)$ with the inside order-one
term $t^2(pE_t+mr\partial_z)$ contributes
\[
 \left(\frac13m'p+mp'-\frac13mm''r,
              \frac23mm'r+m^2r'-mp\right)
\]
to $(A_2,U)$.  Half of its second bracket with the leading term
$trE_t$ contributes
\[
 \left(\frac23mm'r'+\frac12m^2r''+\frac16mm''r,
              -\frac13mm'r-m^2r'\right).
\]
Thus the full effect is
\begin{align}
 \Delta A_2&=\frac13m'p+mp'+\frac23mm'r'
                         +\frac12m^2r''-\frac16mm''r,
 \notag\\
 \Delta U&=-mp+\frac13mm'r.
 \label{jets:second-expansion}
\end{align}
The $m^2r'$ horizontal terms cancel in this sum.

Fix one disk and abbreviate $R=mr$, $m=c/z$, so $r=zR/c$.
Substitution in \eqref{jets:second-expansion} gives
\[
 \begin{aligned}
 \Delta A_2&=-\frac{cR}{z^3}+\frac{cR'}{3z^2}
             +\frac{cR''}{2z}-\frac{cp}{3z^2}+\frac{cp'}z,
 \\
 \Delta U&=-\frac{cp}z-\frac{cR}{3z^2}.
 \end{aligned}
\]
It follows, with every coefficient on the right evaluated at $0$, that
\begin{align}
 F_{ij}:=\Res_0(r_i\Delta A_{2,j}\,dz)
   &=-\frac13R_ip_j-\frac23R_iR_j'-R_i'R_j,\notag\\
 H_{ij}:=\Res_0(r_i'\Delta U_j\,dz)
   &=-R_ip_j-\frac13R_iR_j'-\frac23R_i'R_j.
 \label{jets:laurent-residues}
\end{align}
For example, the double pole $-R_iR_j/z^2$ in the first expression
contributes $-R_i'R_j-R_iR_j'$ to its residue; this accounts for both
derivative terms.  The complete residue of the change in
\eqref{jets:psi} is consequently
\begin{equation}\label{jets:full-residue-matrix}
 \begin{aligned}
 T_{ij}&=F_{ij}+\frac14H_{ij}+\frac34H_{ji}\\
 &=-\frac7{12}R_ip_j-\frac9{12}R_jp_i
                -\frac54R_iR_j'-\frac{17}{12}R_i'R_j.
 \end{aligned}
\end{equation}

Decompose this matrix into half its sum and half its difference with
its transpose.  Directly from \eqref{jets:full-residue-matrix},
\begin{align}
 T_{ij}^{\mathrm{sym}}
 &=-\frac23\bigl(R_ip_j+R_jp_i
                         +2R_iR_j'+2R_i'R_j\bigr),\notag\\
 T_{ij}^{\mathrm{alt}}
 &=\frac1{12}\epsilon_{ij}
                 \bigl((R\wedge p)(0)+(R\wedge R')(0)\bigr).
 \label{jets:symmetric-alternating}
\end{align}
To identify the symmetric term intrinsically, observe that
\[
 P=p+mr'=\frac{R(0)}z+p(0)+2R'(0)+O(z).
\]
Since $P_iP_j$ is a meromorphic section of $L^4=K_C$, its local
expression is $P_iP_j\,dz$, and
\[
 \Res_0(P_iP_j\,dz)
 =R_ip_j+R_jp_i+2R_iR_j'+2R_i'R_j.
\]
Thus the sum of every symmetric entry in
\eqref{jets:symmetric-alternating} is zero, by the residue theorem
on $C$.  For the alternating entry,
\eqref{jets:residue-input} gives $R(0)\wedge p(0)=-1$ and therefore
\[
 \sum_{q\in Z}T_{ij}^{\mathrm{alt}}
 =\frac1{12}\epsilon_{ij}\sum_{\nu=1}^6(W_\nu-1)
 =-\frac34\epsilon_{ij}\deg L.
\]
In particular, the sums of the four entries $(11),(12),(21),(22)$
of the transported inside matrix are respectively
\[
 0,\qquad -\frac34\deg L,\qquad
 \frac34\deg L,\qquad 0.
\]
They coincide with the four outside totals in
\eqref{jets:outside-residue}.  Hence every Serre-duality pairing of the
vertical difference vanishes.  The corresponding class in
$H^1(C,L^3)^2$ is zero, so holomorphic vertical corrections on the
charts complete the order-two matching.
\end{proof}

\subsection{Higher orders and determinant targets}

From order three onward, each step extends both the fields and the
determinant they are required to have.  Write the determinant equation as
\begin{equation}\label{jets:target-equation}
 A\wedge B=t^2G,\qquad G=1+G_1t+G_2t^2+\cdots.
\end{equation}
The field equation in degree $k$ uses $G_{k-2}$.  To compare the next
horizontal rows, this target coefficient must already be compatible
on the overlaps.  We therefore choose each target coefficient before
imposing its field equation, and keep track of these two extensions
separately.

The target is the local bivector $G\sigma^{-1}$.  On a realization $Y$,
with $\cI=\OO_Y(-C)$ and $\cF=\cT_Y(-\log C)$, the relevant line bundle is
\[
 \mathcal M=\cI^4\det\cF=\cI^5\det\cT_Y,
 \qquad \mathcal M|_C=L^4\otimes\cT_C\simeq\OO_C.
\]
Its local frame is $\sigma^{-1}$, and its degree-$k$ associated graded
is $L^k$.  Only its leading trivialization on $C$ is fixed initially;
the successive targets extend that section through finite orders.
After transitions cease to preserve $\sigma$, $G$ is therefore not
an invariant scalar function.  To be explicit, if
$q=\operatorname{div}_{\sigma}D$, then
\begin{equation}\label{jets:target-transport}
 [D,G\sigma^{-1}]=(D(G)-qG)\sigma^{-1},\qquad
 G\longmapsto\exp(D-q)G.
\end{equation}
Here $D-q$ denotes the differential operator $f\mapsto D(f)-qf$.
An order-$k$ transition changes a target with constant term one first
in degree $k$ of $G$, with contribution
$-((k-4)b+h')$.  Once lower coefficients are compatible, a leading
degree-$k$ discrepancy is a cocycle in $L^k$.

The factor $t^2$ in \eqref{jets:target-equation} and the equality
$B_0=0$ determine the order of the choices.  For the field changes
at step $k$, the first determinant equations they can affect are
\[
\begin{array}{c|c|c}
 \text{changed field coefficient}
 &\text{degree in }A\wedge B&\text{target coefficient}\\ \hline
 \text{global addition to }A_{k-1}&k&G_{k-2}\\
 B_k&k&G_{k-2}\\
 A_k&k+1&G_{k-1}.
\end{array}
\]
The target coefficients in the last column remain fixed during these
changes.  In particular, we can still adjust $A_{k-1}$ to solve the
degree-$k$ equation without changing any determinant equation imposed
at an earlier step.  A new target coefficient $G_k$ will be chosen
only after the fields match through index $k$; its field equation
occurs in degree $k+2$.

We will maintain the following invariant after step $n\geq2$:
the field rows through index $n$ match; the determinant equations
\eqref{jets:target-equation} are imposed through degree $n$ in
$A\wedge B$; and the coefficients $G_0=1,G_1,\ldots,G_n$ are chosen
and compatible through degree $n$ under the transition maps.
Thus target compatibility is known two determinant degrees beyond
the equations already imposed: the equations involving $G_{n-1}$
and $G_n$ are still to be solved.
The preceding construction establishes this invariant for $n=2$,
with $G=1$, since both transitions preserve $\sigma$ exactly.

Suppose now that the invariant holds through step $k-1$, where
$k\geq3$.  The coefficient needed in the current determinant equation,
$G_{k-2}$, has already been chosen.  On each chart the horizontal row
must solve
\begin{equation}\label{jets:horizontal-induction}
 r\wedge B_k=G_{k-2}
       -\sum_{a=1}^{k-1}A_a\wedge B_{k-a}.
\end{equation}
This is an equation in $L^{k-2}$.  Indeed, under
\eqref{jets:split-coefficients}, the only additional terms in the
transformation of the sum on the right are multiples of
\[
 \sum_{a=1}^{k-1}B_a\wedge B_{k-a}=0;
\]
the terms cancel in pairs and any middle self-wedge vanishes.

\paragraph{Making the right side divisible at $Z$.}
We may add any row $a\in H^0(C,L^k)^2$ to the previously chosen
vertical coefficient $A_{k-1}$.  Differences of vertical lifts have
this tensor type by \eqref{jets:homogeneous-extension}.
Such a change preserves matching through index $k-1$: every
previous nonsplit transition has positive order, so its action on the
new term $t^kaE_t$ only changes indices at least $k$.
It also preserves all determinant equations already imposed, because
$B_0=0$; its first determinant contribution is $a\wedge B_1$ in
degree $k$.  In particular, it need not preserve the as yet unimposed
equation whose target is $G_{k-2}$.

At each $q_j$ the row $B_1(q_j)=R_j(0)$ is nonzero.
The interpolation in Proposition~\ref{prop:curve-data}(4) gives the
surjection
\[
 H^0(C,L^k)\longrightarrow H^0(Z,L^k|_Z)\qquad(k\geq3).
\]
It permits arbitrary values of both entries of $a$ at all six points.
The scalar maps
$a(q_j)\mapsto a(q_j)\wedge B_1(q_j)$ are surjective, so this
freedom makes the right side of \eqref{jets:horizontal-induction}
vanish at every $q_j$.
Since $r=z\rho$ has only a simple common zero, no derivative
condition is required.  As at order two, division by $z$ on each
disk and Steinness on $U_0$ now give holomorphic solutions $B_k$.

Although this global addition preserves all lower-order matching,
the previous transitions can act on it in the degree currently being solved.
The first transition contributes, at row index $k$,
\begin{equation}\label{jets:old-transition-effect}
 \delta A_k=\frac{k-1}{3}m'a+ma',\qquad
 \delta B_k=-ma.
\end{equation}
All other previous generators, and repeated brackets, affect later
indices.  The displayed horizontal contribution satisfies
$r\wedge(-ma)=a\wedge(mr)$, exactly the corresponding determinant
change on the disk.  We include these terms in the horizontal and
vertical comparisons below.

\paragraph{Matching the horizontal rows.}
Choose any holomorphic vertical lifts of the newly selected $B_k$.
After transport by the previous transitions, the field rows already
agree at every lower index.  The difference of their degree-$k$
determinants is therefore $r\wedge\delta B_k$, since
$\delta A_k\wedge B_0=0$.  Target compatibility through degree
$k-2$ in $G$ makes this difference zero.  Thus the horizontal
discrepancy is a holomorphic multiple of $r$ on each overlap, and
\eqref{jets:leading-action} removes it by a choice of $h$ in an
order-$k$ transition.  Use $b=h'$ when $k=3$, and $b=0$ when
$k\geq4$.
The new transition first affects $G$ in degree $k$, so cannot alter
the target comparison just used in degree $k-2$.

\paragraph{Matching the vertical rows.}
After horizontal matching, the remaining class is in
$H^1(C,L^{k+1})^2$.
For $k\geq4$ it vanishes, since
\begin{equation}\label{jets:vertical-vanishing}
 H^1(C,L^{k+1})\simeq H^0(C,L^{3-k})^*=0.
\end{equation}
Holomorphic vertical corrections on the charts therefore solve the
matching.  They do not change the current determinant equation,
because $B_0=0$.

At $k=3$ the obstruction is $H^1(C,K_C)^2$, tested by the constants
in the two entries.  There is additional overlap freedom to kill it.
On the $j$th overlap replace $h$ by $h+\lambda_j/z$ and compensate
on its disk by
\begin{equation}\label{jets:third-compensation}
 B_3\longmapsto B_3+\frac{3\lambda_jr}{z}.
\end{equation}
This is holomorphic and its wedge with $r$ is zero.  It preserves
both the current determinant equation and horizontal matching:
the transition adds the opposite horizontal row $-3\lambda_jr/z$.
Keep the vertical lift $A_3$ unchanged in the chosen disk coordinate.
This prescribes a holomorphic section of $\mathcal E_3$ on that disk;
in another split coordinate its vertical component acquires the term
$a\,\delta B_3$ dictated by \eqref{jets:split-coefficients}.
Changing a degree-three input row affects its transport under old
positive-order transitions only at later indices.
The change in the transported vertical difference at index three is
therefore exactly
\begin{equation}\label{jets:third-residue}
 \lambda_j\left(\frac{r'}z+\frac{2r}{z^2}\right),\qquad
 \Res_0\left[\lambda_j\left(\frac{r'}z+\frac{2r}{z^2}\right)dz\right]
       =3\lambda_jr'(0).
\end{equation}
The two entries here are coefficients of canonical differentials,
since $L^4=K_C$.  The six rows $r'(q_j)$ span $\CC^2$: their
projective directions are those of the pencil values $u(q_j)$, which
are not all the same.  Hence the map
$(\lambda_1,\ldots,\lambda_6)\mapsto\sum_j3\lambda_jr'(q_j)$
has rank two and can cancel both obstruction entries.
The resulting vertical class splits, completing order three.

The correction \eqref{jets:third-compensation} can change the
degree-four coefficient of $A\wedge B$.  This causes no conflict:
$G_2=0$ has been selected, but its determinant equation is first
imposed at step four.  The degree-three equation already imposed is
unchanged.  Moreover the modified order-three transition still has
$b=h'$ and so preserves $\sigma$ exactly, leaving the transformation
law of the chosen targets intact.

\paragraph{Extending the determinant target.}
It remains to choose $G_k$ compatibly after the order-$k$ transition.
For $k=3$ all transitions preserve $\sigma$, so we retain $G=1$.
At $k=4$ the earlier transitions still preserve it exactly; thus the
first target discrepancy is $-h'$ in degree four, by
\eqref{jets:target-transport} with $b=0$.
The horizontal coefficient of an order-four generator has type
$L^4\otimes\cT_C=\OO_C$.  More explicitly, its split
transformation is $\widetilde h=g^{4-k}h$, which equals $h$ at
$k=4$.  Consequently $h'dz=dh$ is an intrinsic exact differential
on the overlap.  Each such differential has zero contour residue.
The sole obstruction to splitting this $L^4=K_C$ cocycle is its
pairing with the constant section of $\OO_C$, so it vanishes.
There are holomorphic chart coefficients $G_4$ accounting for the
discrepancy, with $G_1=G_2=G_3=0$.
Transport by a previous positive-order modification changes a
degree-four term only in degrees at least five. Equivalently, conjugating
the order-four generator by the earlier transitions leaves its degree-four
part unchanged. Thus the scalar transformation law of $h$, and the residue
calculation using $dh$, apply intrinsically at this order.

For $k>4$, once the lower target coefficients are compatible, the
new leading discrepancy lies in $L^k$, and
\begin{equation}\label{jets:target-vanishing}
 H^1(C,L^k)\simeq H^0(C,L^{4-k})^*=0.
\end{equation}
Thus it always splits and defines $G_k$.
In every case this target extension leaves all equations already
imposed unchanged, since $G_k$ is used only in degree $k+2$ of
$A\wedge B$.  This completes the induction invariant.

We now choose where to stop.  A row of index $k$ contributes degree
$k+1$ in the logarithmic frame $(E_t,\partial_z)$, so matching through
index seven prescribes a section modulo $\cI^9\cF$.  This is the
threshold needed in Section~\ref{analytic:section}: the successive
quotients of the unprescribed tail are extensions of $L^{n-4}$ by
$L^n$, and both line bundles have vanishing first cohomology for
$n\geq9$.  Lemma~\ref{analytic:actual-lifting} turns these graded
vanishings into holomorphic lifting on a realized neighborhood.

\begin{proposition}[Finite compatible jets]\label{prop:jets}
For the data of Propositions~\ref{prop:curve-data} and
\ref{prop:residue-data}, there are finitely many holomorphic
disk-to-outside transition germs, each a composition of flows
\eqref{jets:generator} of orders $1,\ldots,7$, together with
holomorphic local rows $A_k,B_k$ for $0\leq k\leq7$, such that:
the transitions fix $C$ and preserve its conormal bundle $L$;
the vector fields \eqref{jets:fields} match through row index seven;
and $A\wedge B=t^2(1+O(t))$ on every chart.

On any surface neighborhood $Y$ realizing these transition germs,
write $\cI=\OO_Y(-C)$ and $\cF=\cT_Y(-\log C)$.
The rows give two sections of $\cF/\cI^9\cF$, lying in
$\cI\cF/\cI^9\cF$.  Their determinant has leading term
$t^4E_t\wedge\partial_z$, equivalently
$t^5\partial_t\wedge\partial_z$, with coefficient one.
\end{proposition}

\begin{proof}
Perform the first two steps and then the induction through $k=7$.
Only finitely many flows occur, and each is an actual holomorphic
germ on an annular neighborhood of the zero section, as explained
after Lemma~\ref{jets:bracket}.  Terms with row indices $0,\ldots,7$
have logarithmic-frame coefficients of degrees $1,\ldots,8$.
Matching those terms is precisely matching modulo $\cI^9\cF$.
Finally \eqref{jets:determinant} and the degree-two determinant
equation give
\[
 v_1\wedge v_2=t^4(1+O(t))E_t\wedge\partial_z
             =t^5(1+O(t))\partial_t\wedge\partial_z.
\]
This statement uses only finite compatible jets; analytic realization
and the lifting of these quotient sections are addressed next.
\end{proof}
